\section{The sign component of the Theta conjecture at \texorpdfstring{$t=0$}{t=0}}
\label{sec:theta-sign}

We prove the sign-isotypic component of the $t=0$ specialization of
the Theta coinvariant conjecture of D'Adderio, Iraci and Vanden
Wyngaerd \cite[Conjecture~8.2]{DIV2021}, in every rank.  In particular,
the alternating coinvariants with one bosonic and two fermionic
alphabets have dimension $3^{n-1}$, proving the numerical prediction
in Bergeron's Table~3 \cite{BergeronSuper2020}.
Iraci, Nadeau and Vanden Wyngaerd
\cite[Theorem~5.6]{INV2024} computed the symmetric-function prediction;
the argument here identifies that prediction with the actual
coinvariant sign component.  Its ingredients are the full hook
identity, diagonal supersymmetry, and the Fields character theorem.

\subsection{Adjoining a fermionic alphabet}

For $k,j\geq0$, let
\begin{align*}
 P_n^{(k,j)}
   &=K[x_i^{(a)}:1\leq i\leq n,\ 1\leq a\leq k]
       \otimes\Lambda(\theta_i^{(b)}:1\leq i\leq n,\ 1\leq b\leq j),\\
 R_n^{(k,j)}
   &=P_n^{(k,j)}/\langle(P_n^{(k,j)})^{S_n}_+\rangle.
\end{align*}
Here $K$ has characteristic zero, all odd generators anticommute,
and $S_n$ acts by simultaneously permuting their label indices.
The subscript $+$ denotes zero constant term.  We record each
bosonic and fermionic alphabet separately, by
$Q=(q_1,\ldots,q_k)$ and $U=(u_1,\ldots,u_j)$, respectively.
Supersymmetric evaluation $f[Q\mid U]$ uses the positive convention
of \eqref{eq:superschur}; in particular, odd degrees have positive
coefficients in the Hilbert series.

Write $W_n=\ker(K^n\xrightarrow{\sum}K)$.  If
$\eta_1,\ldots,\eta_n$ is a new fermionic alphabet, its invariant
$\sum_i\eta_i$ is zero in $R_n^{(k,j+1)}$.  The quotient
$K^n/K\mathbf1$ identifies canonically with $W_n$ by
$z\mapsto z-n^{-1}(\sum_i z_i)\mathbf1$.
Consequently the quotient map factors through
$R_n^{(k,j)}\otimes\Lambda W_n$; the tensor product uses the Koszul
sign rule with the old odd variables.

\begin{theorem}[Natural sign transfer in superspace]
\label{thm:all-superspace-sign-transfer}
For all $n\geq1$ and $k,j\geq0$, the natural quotient induces an
isomorphism
\begin{equation}\label{eq:all-superspace-sign-transfer}
 \left(R_n^{(k,j)}\otimes\Lambda W_n\right)^{\sgn}
 \xrightarrow{\ \sim\ }
 \left(R_n^{(k,j+1)}\right)^{\sgn}.
\end{equation}
It preserves every old grading and the degree in the new fermionic
alphabet, and commutes with parity-preserving changes of the old
row alphabets.  In new fermionic degree $0\leq d\leq n-1$, the
induced isomorphism of multiplicity spaces is
\begin{equation}\label{eq:super-hook-multiplicity}
 \Hom_{S_n}\!\left(S^{(d+1,1^{n-d-1})},R_n^{(k,j)}\right)
 \cong
 \Hom_{S_n}\!\left(\sgn,
       (R_n^{(k,j+1)})_{\deg_\eta=d}\right).
\end{equation}
Both sides of \eqref{eq:all-superspace-sign-transfer} vanish in
new fermionic degree greater than $n-1$.
\end{theorem}

The case $(k,j)=(2,0)$ is the sign-transfer theorem of
Gonz\'alez, Lentfer and Mularczyk \cite[Theorem~4.21]{GLM}.
Here the full hook identity permits arbitrary old bosonic and
fermionic alphabets.

\begin{proof}
The quotient is surjective on sign components, since the sign
idempotent is exact in characteristic zero.  By Lentfer's diagonal
supersymmetry theorem \cite[Theorem~1.2 and Corollary~1.3]{Lentfer},
the character of the $S^\mu$ multiplicity space of $R_n^{(k,j)}$ is
$C_\mu^{(n)}[Q\mid U]$.  Theorem~\ref{thm:full-hook} gives
\[
 C_{(d+1,1^{n-d-1})}^{(n)}=e_d^\perp A_n
 \qquad(0\leq d\leq n-1).
\]
Super Schur branching in one additional odd variable $v$ yields
\begin{equation}\label{eq:super-sign-branching}
 A_n[Q\mid U,v]
   =\sum_{d\geq0}v^d(e_d^\perp A_n)[Q\mid U].
\end{equation}
The skew operator acts on the stable symmetric function before
the finite supersymmetric specialization is made.

The invariant dot product on $W_n$ identifies its exterior powers
with their duals, and
$\sgn\otimes\Lambda^dW_n\cong S^{(d+1,1^{n-d-1})}$.
Tensor--Hom adjunction therefore identifies the coefficient of
$v^d$ in the source sign multiplicity with
$C_{(d+1,1^{n-d-1})}^{(n)}[Q\mid U]$.
This equals the coefficient in \eqref{eq:super-sign-branching},
which is the target sign multiplicity.  Every multigraded piece is
finite-dimensional, so the surjection is an isomorphism in every
multidegree, proving both displayed isomorphisms.  The source has
no exterior degree above $n-1$; the target also has none, because
the relation $\sum_i\eta_i=0$ reduces its new exterior alphabet to
$W_n$.  Naturality follows from the quotient map itself.
\end{proof}

We use the usual grading convention when passing to harmonic
models: a harmonic polynomial and its coinvariant class have the
same multidegree.  More explicitly, fix an order on the odd
generators and use the apolar pairing
\[
 \langle x^\alpha\theta_I,x^\beta\theta_J\rangle
   =\alpha!\,\delta_{\alpha,\beta}\delta_{I,J}.
\]
The odd differential operators are taken in the reverse order to
the corresponding monomial.  The harmonic inverse system is the
orthogonal complement of the invariant ideal for this pairing.
The pairing is positive definite over the reals and is defined
over $\Q$, so the quotient projection identifies the harmonic
model and the coinvariant module in each multidegree over every
characteristic-zero field.  Thus no inversion of $q$, $u$, or $v$
occurs when using the Fields formula \eqref{eq:fields-character}.
On harmonic inverse systems, \eqref{eq:all-superspace-sign-transfer}
also gives
\[
 (\Harm_n^{(k,j+1)})^{\sgn}
   =(\Harm_n^{(k,j)}\otimes\Lambda W_n)^{\sgn}:
\]
the final harmonic space is contained in the intermediate one,
and the quotient isomorphism makes their sign components equal.

\subsection{The actual alternating character}

Put $[m]_q!=\prod_{r=1}^m(1+q+\cdots+q^{r-1})$, with
$[0]_q!=1$, and use the Gaussian coefficients
\[
 \genfrac{[}{]}{0pt}{}{N}{a,b,c}_q
    =\frac{[N]_q!}{[a]_q![b]_q![c]_q!}
 \quad(a+b+c=N),
 \qquad
 \genfrac{[}{]}{0pt}{}{N}{a}_q
    =\frac{[N]_q!}{[a]_q![N-a]_q!}.
\]

\begin{theorem}[The $t=0$ Theta sign component]
\label{thm:one-boson-two-fermion-sign}
For every $n\geq1$,
\begin{equation}\label{eq:one-boson-two-fermion-sign}
 \operatorname{Hilb}\!\left((R_n^{(1,2)})^{\sgn};q;u,v\right)
 =\sum_{\substack{a,b\geq0\\a+b\leq n-1}}
 u^av^b q^{\binom{n-a-b}{2}}
 \genfrac{[}{]}{0pt}{}{n-1}{a,b,n-1-a-b}_q.
\end{equation}
Here $q$ records the bosonic degree, and $u,v$ record the two
fermionic degrees.  This is the sign coefficient of
\cite[Conjecture~8.2]{DIV2021} at $t=0$, in every fermionic
bidegree.\footnote{Equation~(76) in the arXiv version of
\cite{DIV2021} prints $1\leq a+b<n$.  We include the ordinary
term $a=b=0$, whose sign character $q^{\binom n2}$ is classical.}
In particular,
\begin{equation}\label{eq:one-boson-two-fermion-sign-dimension}
 \operatorname{Hilb}\!\left((R_n^{(1,2)})^{\sgn};1;u,v\right)
 =(1+u+v)^{n-1},
 \qquad
 \dim(R_n^{(1,2)})^{\sgn}=3^{n-1}.
\end{equation}
\end{theorem}

\begin{proof}
Fix $0\leq b\leq n-1$ and set $N=n-1$, $d=N-b$.
Theorem~\ref{thm:all-superspace-sign-transfer} identifies the
coefficient of $v^b$ with the multiplicity of the hook
$\mu=(b+1,1^d)$ in $R_n^{(1,1)}$.
The Fields character theorem of Murai, Rhoades and Wilson
\cite[Theorem~5.2 and equation~(2.5)]{MRW}, in the form
\eqref{eq:fields-character}, evaluates this multiplicity as
\[
 \sum_{T\in\SYT(\mu)}q^{\maj(T)}
                  \prod_{r=1}^{\des(T)}(1+u q^{-r}).
\]
The hook-tableau enumeration used in \eqref{eq:hook-hook-formula}
gives $\operatorname{des}(T)=d$ and
$\sum_Tq^{\maj(T)}=q^{\binom{d+1}{2}}
 \genfrac{[}{]}{0pt}{}{N}{d}_q$.
It follows that
\begin{equation}\label{eq:theta-sign-fields-product}
 [v^b]\operatorname{Hilb}\!\left((R_n^{(1,2)})^{\sgn};q;u,v\right)
   =q^{\binom{d+1}{2}}\genfrac{[}{]}{0pt}{}{N}{b}_q
                \prod_{r=1}^{d}(1+u q^{-r}).
\end{equation}
For $0\leq a\leq d$, the Gaussian binomial theorem gives
\[
 [u^a]\prod_{r=1}^{d}(1+u q^{-r})
   =\sum_{\substack{J\subseteq[d]\\|J|=a}}q^{-\sum_{r\in J}r}
   =q^{-ad+\binom a2}\genfrac{[}{]}{0pt}{}{d}{a}_q.
\]
For $a>d$ the coefficient is zero.  Since
\[
 \binom{d+1}{2}-ad+\binom a2=\binom{d-a+1}{2}
      =\binom{n-a-b}{2},
 \qquad
 \genfrac{[}{]}{0pt}{}{N}{b}_q
 \genfrac{[}{]}{0pt}{}{N-b}{a}_q
     =\genfrac{[}{]}{0pt}{}{N}{a,b,N-a-b}_q,
\]
extracting $u^a$ from \eqref{eq:theta-sign-fields-product} proves
\eqref{eq:one-boson-two-fermion-sign}.

For comparison with the stated conjecture, let
$\Theta_f$ and $\nabla$ denote the symmetric-function operators
of \cite{DIV2021}.  The sign coefficient predicted there in
fermionic bidegree $(a,b)$, after setting $t=0$, is
\[
 \left\langle
     \left.\Theta_{e_a}\Theta_{e_b}\nabla e_{n-a-b}\right|_{t=0},
     e_n\right\rangle
   =q^{\binom{n-a-b}{2}}
      \genfrac{[}{]}{0pt}{}{n-1}{a+b}_q
      \genfrac{[}{]}{0pt}{}{a+b}{a}_q
 \qquad(a+b<n),
\]
by \cite[Theorem~5.6]{INV2024}; see also
\cite[Theorem~8.3]{Lentfer12}.  The two Gaussian binomials multiply
to the multinomial in \eqref{eq:one-boson-two-fermion-sign}.
Both series have zero sign coefficient for $a+b\geq n$.
Finally, putting $q=1$ in that equation and applying the
trinomial theorem proves
\eqref{eq:one-boson-two-fermion-sign-dimension}.
\end{proof}
